\section{Method and Worked Example}
% Define the symbols, explain the pipeline, and show the worked example.
\subsection{Method}

\subsubsection{LoRA}

\paragraph{Single-task update}

For a pre-trained weight matrix \(W_0\in\mathbb{R}^{d\times k}\), LoRA constrains its update to \(\Delta W=AB\), where \(A\in\mathbb{R}^{d\times r}\), \(B\in\mathbb{R}^{r\times k}\), and \(r\ll\min(d,k)\). With input \(x\in\mathbb{R}^{k}\), the adapted layer computes
\[
	h=W_0x+\frac{\alpha}{r}ABx
\]
where \(\alpha\) is a scaling hyperparameter.
The factor \(\alpha/r\) reduces the need to retune other hyperparameters when the rank \(r\) changes.
Only \(A\) and \(B\) are trained; \(W_0\) remains frozen.

For this matrix, LoRA trains \(r(d+k)\) parameters instead of the \(dk\) weights that full fine-tuning would update. This reduces the task-specific checkpoint size and the gradient and optimizer-state memory for the adapted matrix. The base weights must still be stored and used.


\paragraph{Training and inference}

The original LoRA method initializes \(A\) with random Gaussian values and \(B\) with zeros. Consequently, \(AB=0\) at the start of training, so the adapted layer initially computes exactly \(W_0x\). Training then changes \(A\) and \(B\) while leaving \(W_0\) fixed.

After training, the update can be merged into the weight matrix:
\[
	W_{\mathrm{merged}}=W_0+\frac{\alpha}{r}AB
\]
Inference with \(W_{\mathrm{merged}}x\) uses the same matrix multiplication as the original layer and incurs no additional latency from a separate LoRA path. For multiple tasks, one can retain the shared base model and store a small adapter for each task, merging the chosen adapter when needed. Deploying an adapter separately still requires access to its matching base model.

\subsubsection(O-LoRA)

\paragraph{Task-specific LoRA updates}

Suppose tasks \(\mathcal{D}_1,\ldots,\mathcal{D}_T\) arrive sequentially, and only the current task's data \(\mathcal{D}_t\) are available when learning task \(t\). For a pre-trained weight matrix \(W_0 \in \mathbb{R}^{d \times k}\), O-LoRA learns a separate low-rank update \(A_t B_t\) for each task, where \(A_t \in \mathbb{R}^{d \times r}\), \(B_t \in \mathbb{R}^{r\times k}\), and \(r \ll \min(d,k)\). After task \(t\), the adapted matrix is
\[
	W_t = W_0 + \sum_{i=1}^{t} \Delta W_i
\]
where \(\Delta W_i = \frac{\alpha}{r} A_i B_i\).
With input \(x\in\mathbb{R}^{k}\), the adapted layer computes
\[
	h_t = W_t x
\]
The base matrix and the previous task factors \(\{A_i,B_i\}_{i<t}\) remain fixed while \(A_t\) and \(B_t\) are trained.
The complete model update \(\Delta\Phi_i\) collects these updates across all selected weight matrices. Thus,
\[
	\Phi_t = \Phi_0+\sum_{i=1}^{t}\Delta\Phi_i
\]
Here, \(\Phi_t\) denotes the complete model parameters after incorporating the updates from tasks \(1,\ldots,t\), while \(W_t\) denotes an individual adapted weight matrix.
For brevity, the matrix-level formulation is shown for one selected weight matrix; the same updates and orthogonality penalties are applied to all selected matrices.
Then the task loss is
\[
	\mathcal{L}_{\mathrm{task}} =
	- \sum_{(x,y)\in\mathcal{D}_t}\log P_{\Phi_t}(y\mid x)
\]

\paragraph{Orthogonality constraint}

O-LoRA approximates the task-specific update subspace \(\mathcal U_t \subseteq \mathbb{R}^{d}\) by the column space of \(A_t\) for task \(t\). To make \(\mathcal{U}_t\) orthogonal to a previous task's subspace \(\mathcal{U}_i\), O-LoRA would require \(A_i^{\top}A_t=0\). It encourages this condition with the penalty
\[
	\mathcal{L}_{\mathrm{orth}}(A_t)
	= \sum_{i=1}^{t-1} \left\|A_i^{\top} A_t\right\|_F^2
	= \sum_{i=1}^{t-1} \sum_{p=1}^{r}\sum_{q=1}^{r} \left((A_i^{\top} A_t)_{pq}\right)^2
\]
where \(\|\cdot\|_F\) denotes the Frobenius norm of a matrix.

For a positive weight \(\lambda\), the current task's training objective can be written as
\[
	\mathcal{L}_t =
	\mathcal{L}_{\mathrm{task}}
	+ \lambda \mathcal{L}_{\mathrm{orth}}(A_t)
\]
The model fits the current task while penalizing overlap with previous update subspaces. A finite penalty encourages orthogonality but does not guarantee it.

\paragraph{Training and inference}

In the paper, LoRA is applied to the attention query and value weights. Training task \(t\) needs the current data and the frozen earlier LoRA factors, but no stored examples from previous tasks. The number of task factors grows during training. For inference, their updates can be merged into the pre-trained weights as in \(W_t\), so the model can make predictions without selecting an adapter using a task ID.
